% ============================================================
\section{Validation Evidence}
% ============================================================

\subsection{Geometric accuracy: in-vitro concurrent validity}

Masch et al.\ evaluated the intrinsic measurement accuracy of the sensor using a
standardised in-vitro protocol to eliminate the confounding influence of soft-tissue
artefacts and inter-individual anatomical variation \cite{walkling2025}.
Ten sensors were each applied ten times to precision 3D-printed templates representing
defined flexion radii and specific torsion angles.

Across all flexion templates, the \ac{MAE} was:

\begin{itemize}
	\item \SI{1.2}{\milli\metre} for 3D segment position.
	\item \ang{1.1} for overall angular orientation.
\end{itemize}

On the spine-shaped template (anatomically realistic curvature):

\begin{itemize}
	\item Position \acs{MAE}: \SI{0.7}{\milli\metre}.
	\item Angular \acs{MAE}: \ang{0.8}.
	\item Torsion \acs{MAE}: \ang{0.4}.
\end{itemize}

\Cref{tab:accuracy} places these values in the context of other systems.

\begin{table}[ht]
	\centering
	\caption{Comparative measurement accuracy across spine assessment systems.}
	\label{tab:accuracy}
	\small
	\begin{tabularx}{\textwidth}{%
		>{\bfseries\raggedright\arraybackslash}p{42mm}%
		>{\raggedright\arraybackslash}p{27mm}%
		>{\raggedright\arraybackslash}X%
		>{\raggedright\arraybackslash}X}
		\toprule
		System                                          & Primary use             & Angular error & Measurement mode \\
		\midrule
		\rowcolor{rowgray}
		Optical \acs{MoCap} (Vicon)                     & Lab biomechanics        &
		\SI{0.154}{\milli\metre} RMSE~\cite{vicon_rmse} &
		Static \& dynamic (lab only)                                                                                 \\
		Idiag M360 (SpinalMouse)                        & Clinical diagnostics    &
		\ang{1}--\ang{2}~\cite{spinalMouse}             &
		Static (manual sweep)                                                                                        \\
		\rowcolor{rowgray}
		FlexTail (this system)                          & Research \& ambulatory  &
		\ang{0.8}--\ang{1.1}~\cite{walkling2025}        &
		Static \& dynamic (continuous)                                                                               \\
		BodyGuard (occupational wearable)               & Occupational ergonomics &
		$\approx\ang{2.5}$                              &
		Dynamic (inclination only)                                                                                   \\
		\rowcolor{rowgray}
		Standard multi-IMU array                        & Field studies           &
		$>\ang{5}$ incl.\ yaw drift~\cite{madgwick2011} &
		Dynamic (drift-limited)                                                                                      \\
		\bottomrule
	\end{tabularx}
\end{table}

The FlexTail error range matches clinical-grade static instruments while enabling
continuous ambulatory use --- a combination not demonstrated by the comparison
technologies reviewed here.

\begin{figure}[!b]
	\centering
	\begin{tikzpicture}
		\begin{axis}[
				xbar,
				width=0.82\textwidth,
				height=6.5cm,
				bar width=10pt,
				xmin=0, xmax=6.5,
				xlabel={Angular \acs{MAE} (\textdegree)},
				symbolic y coords={
						Multi-IMU with drift,
						BodyGuard,
						SpinalMouse,
						FlexTail all templates,
						FlexTail spine template
					},
				ytick=data,
				nodes near coords,
				nodes near coords align={horizontal},
				enlarge y limits=0.18,
				axis line style={draw=gray!60},
				tick style={draw=gray!60},
				grid=major,
				grid style={dashed, gray!30},
				extra x ticks={2},
				extra x tick labels={},
				extra x tick style={
						grid=major,
						grid style={red!60!black, dashed, thick},
						tick style={draw=none},
					},
			]
			\addplot[
				fill=minktecBlue!80, draw=minktecBlue!80
			] coordinates {
					(5.5,Multi-IMU with drift)
					(2.5,BodyGuard)
					(1.5,SpinalMouse)
					(1.1,FlexTail all templates)
					(0.8,FlexTail spine template)
				};
		\end{axis}
	\end{tikzpicture}
	\caption{Angular mean absolute error comparison across spine assessment systems. The dashed red line at \ang{2} marks an informal clinical-use threshold. FlexTail on the spine-shaped template achieves \ang{0.8} MAE, placing it below all non-optical competitors. Values from \cite{walkling2025, spinalMouse, madgwick2011}; BodyGuard value from manufacturer specification.}
	\label{fig:accuracy_bar}
\end{figure}

\subsection{Human activity recognition: comparison with video-based pose estimation}

Walkling et al.\ (2025) conducted a direct comparison between FlexTail and a
camera-based pose estimation model --- specifically OpenPose~1.7.0, which extracts
25 body joint keypoints per frame --- across eleven \acp{ADL} \cite{walkling2025}.
Ten healthy participants performed the protocol indoors, wearing the sensor shirt
while a smartphone camera mounted on a tripod captured each session simultaneously.

The authors selected the \ac{RDST} classifier for FlexTail time-series data and the
QUANT classifier for the OpenPose keypoint sequences, both chosen as the best-performing
algorithm for their respective modality from a benchmark of six state-of-the-art
time-series classifiers.
The dataset is small (10 subjects, $\approx$100 min total), and the authors
explicitly note that performance differences between the two systems cannot be tested
for statistical significance at this sample size.

Key results:

\begin{itemize}
	\item At a \textbf{1-second} classification window, both systems achieved a mean
	      \textbf{F1-score of 0.90} --- RDST for FlexTail and QUANT for camera data.
	\item For \textbf{standing-up transitions}, the wearable achieved \textbf{95\,\%}
	      accuracy versus \textbf{76\,\%} for the OpenPose model.
\end{itemize}

The camera model outperformed the wearable on activities that produce large, visually
distinctive arm and hand movements --- eating, food cutting, and dishwasher
loading/unloading --- confirming that camera-based and spine-based sensing are
complementary rather than mutually exclusive.
For trunk-posture and transition categories most relevant to clinical and ergonomic
monitoring, the wearable provided consistently higher accuracy.

The causal explanation is straightforward: the strain-gauge array provides an
uninterrupted, occlusion-free kinematic signal of the spinal column.
A pose estimation model inferring trunk kinematics from external joint-keypoint
coordinates is sensitive to viewpoint variation and leg-crossing artefacts --- exactly
the conditions that arise when the camera position is varied between participants, as
it was in this study.

\subsection{Post-surgical range of motion: clinical pilot}

Marx et al.\ (2023) used the FlexTail sensor in a biomechanical pilot to measure \ac{ROM}
before and after posterior corrective spinal fusion surgery (dorsale
Korrekturspondylodese) \cite{marx2023}.
The protocol quantified maximum segmental flexion, extension, lateral flexion, and
torsion pre- and post-operatively, providing a continuous, objective complement to
patient-reported outcome measures.
The authors report accurate real-time 3D representation as the key methodological
advantage over static radiological assessment, which captures morphology but not
functional kinematics under load.
